% GENERATED FILE. Edit canonical lessons, then run npm run generate:books.
% canonical-source-hash: fnv1a64:10ca3ef890ed7182
% canonical-lessons: ZH-C159-woshi, ZH-C159-keting, ZH-C159-shufang, ZH-W159-fang-room, ZH-C159-yangtai, ZH-C159-chuanghu

\chapter{Rooms at Home}
\label{ch:zh-rooms-at-home}
\hlchaptermodality{159}

\begin{chapteropening}
\textbf{By the end of this chapter:} \emph{I can name the bedroom, the living room, the study, the balcony and the window.}

Complete the last lesson of chapter 159: I can name the bedroom, the living room, the study, the balcony and the window.
\end{chapteropening}

\section[wòshì]{wòshì — bedroom}

\label{lesson:ZH-C159-woshi}



\begin{quote}
Before the new one: say the Mandarin for to stand up, then the Mandarin for to be used to.
\end{quote}

\subsection*{You'll want to know: wòshì}
\emph{wòshì} — \textquotedblleft{}bedroom\textquotedblright{}.

\begin{quote}
\emph{wǒ de wòshì hěn xiǎo} — My bedroom is small.
\end{quote}

Its characters come later; for now it is a word for the ear.

\begin{cousinweb}
\emph{wòshì} is built from \emph{wò} 'to lie down' and \emph{shì} 'room', both inherited from Middle Chinese, \emph{ngwaH} and \emph{syit} (Guangyun, in Baxter's transcription): a room to lie down in.
\end{cousinweb}

\begin{grammarlens}[title={What you've built}]
One more word for everyday things.
\end{grammarlens}

\subsection*{Your turn}
\begin{itemize}
\raggedright
  \item \emph{Say it:} \emph{wòshì}
  \item \emph{Say it:} \emph{wòshì}, once more
  \item \emph{Say it:} \emph{wǒ de wòshì hěn xiǎo} — My bedroom is small.
  \item \emph{Recall:} say the Mandarin for to stand up, then the Mandarin for to be used to, then say \emph{wòshì} again
\end{itemize}

\begin{tcolorbox}[breakable,colback=teal!4,colframe=teal!35!black,title={Before you move on}]
What does wòshì mean? (\textquotedblleft{}Bedroom\textquotedblright{}.) Say it once more.
\end{tcolorbox}
\section[kètīng]{kètīng — living room; sitting room}

\label{lesson:ZH-C159-keting}



\begin{quote}
Before the new one: say the Mandarin for to be used to, then the Mandarin for a bedroom.
\end{quote}

\subsection*{You'll want to know: kètīng}
\emph{kètīng} — \textquotedblleft{}living room; sitting room\textquotedblright{}.

\begin{quote}
\emph{kèrén zài kètīng} — The guests are in the living room.
\end{quote}

Its characters come later; for now it is a word for the ear.

\begin{cousinweb}
\emph{kètīng} is built from \zh{客} 'guest' and \emph{tīng} 'hall'. The first is inherited from Middle Chinese \emph{khaek} (Guangyun, in Baxter's transcription); the second is listed there as \emph{theng}, glossed 'a hall', a different word from the \emph{tīng} of listening: the hall for guests.
\end{cousinweb}

\begin{grammarlens}[title={What you've built}]
One more word for everyday things.
\end{grammarlens}

\subsection*{Your turn}
\begin{itemize}
\raggedright
  \item \emph{Say it:} \emph{kètīng}
  \item \emph{Say it:} \emph{kètīng}, once more
  \item \emph{Say it:} \emph{kèrén zài kètīng} — The guests are in the living room.
  \item \emph{Recall:} say the Mandarin for to be used to, then the Mandarin for a bedroom, then say \emph{kètīng} again
\end{itemize}

\begin{tcolorbox}[breakable,colback=teal!4,colframe=teal!35!black,title={Before you move on}]
What does kètīng mean? (\textquotedblleft{}Living room; sitting room\textquotedblright{}.) Say it once more.
\end{tcolorbox}
\section[\texorpdfstring{\zh{书房} (shūfáng)}{书房 (shūfáng)}]{\zh{书房} (shūfáng) — study; a room for books and work}

\label{lesson:ZH-C159-shufang}



\begin{quote}
Before the new one: say the Mandarin for a bedroom, then the Mandarin for a living room.
\end{quote}

\subsection*{You'll want to know: \zh{书房}}
\textbf{\zh{书房}} — \emph{shūfáng} — \textquotedblleft{}study; a room for books and work\textquotedblright{}.

\begin{quote}
\textbf{\zh{我的书都在书房里}} \emph{wǒ de shū dōu zài shūfáng lǐ} — All my books are in the study.
\end{quote}

The new character, in pieces: \zh{房} — eight strokes in two pieces: on top a door of four strokes (a dot, a horizontal that turns down, a horizontal closing it, a long stroke falling to the left), and under it the four strokes of the piece read fāng. The next lesson puts it in your hand.

\begin{cousinweb}
\textbf{\zh{书房}} (\emph{shūfáng}) is built from \zh{书} 'book' and \zh{房} 'room', both inherited from Middle Chinese, \emph{syo} and \emph{bjang} (Guangyun, in Baxter's transcription), which glosses the second 'a room': the room for books.
\end{cousinweb}

\begin{grammarlens}[title={What you've built}]
One more word for everyday things.
\end{grammarlens}

\subsection*{Your turn}
\begin{itemize}
\raggedright
  \item \emph{Say it:} \emph{shūfáng}
  \item \emph{Say it:} \emph{shūfáng}, once more
  \item \emph{Say it:} \emph{wǒ de shū dōu zài shūfáng lǐ} — All my books are in the study.
  \item \emph{Recall:} say the Mandarin for a bedroom, then the Mandarin for a living room, then say \emph{shūfáng} again
\end{itemize}

\begin{tcolorbox}[breakable,colback=teal!4,colframe=teal!35!black,title={Before you move on}]
What does \zh{书房} mean? (\textquotedblleft{}Study; a room for books and work\textquotedblright{}.) Say it once more.
\end{tcolorbox}
\section[\texorpdfstring{\zh{房} (fáng)}{房 (fáng)}]{\zh{房} — a door over fāng}

\label{lesson:ZH-W159-fang-room}



\begin{quote}
The word first, then the character inside it.

\begin{itemize}
\raggedright
  \item \emph{Recall:} say \emph{shūfáng} — \textbf{R1}, one lesson back, before it decays
\end{itemize}

\begin{quote}
\zh{书} \zh{门}
\end{quote}

\emph{Read it:} each one aloud, with nothing said first — \emph{shū}, \emph{mén}
\end{quote}

\subsection*{Script — a door over fāng}
\hlblockfigure{\detokenize{figures/ZH-W159-fang-room-filmstrip.pdf}}{How \zh{房} is written, stroke by stroke}

\begin{quote}
\zh{房}
\end{quote}

\textbf{Two pieces, top over bottom.} Count the pieces, not the strokes.

\textbf{The door, four strokes.} A dot on top. A horizontal that turns down at its right end. A horizontal under it, closing the box. Then a long stroke from the top left falling down to the lower left: the side of the door.

\textbf{The piece read fāng, four strokes, under the door.} A dot. A long horizontal. A horizontal that turns down and hooks back to the left. Last, a stroke falling to the left through the middle.

4 + 4 = eight.

\subsection*{Your turn}
\begin{itemize}
\raggedright
  \item \emph{Write it:} \zh{房}, all eight, with the model in front of you
  \item \emph{Cover:} write it again from memory, and count the pieces, not the strokes
\end{itemize}

\begin{tcolorbox}[breakable,colback=teal!4,colframe=teal!35!black,title={Before you move on}]
How many strokes in \zh{房}? (\textbf{Eight}.) Which piece is on top? (\textbf{The door, four strokes}.) Which stroke of the lower piece hooks? (\textbf{The third, the one that turns down}.)

Source: \href{https://raw.githubusercontent.com/chanind/hanzi-writer-data/68d10a4b21150cae5e1ebbd223eed289cf32d90c/data/\%E6\%88\%BF.json}{Hanzi Writer Data \zh{房}, pinned stroke paths and medians}.
\end{tcolorbox}
\section[yángtái]{yángtái — balcony}

\label{lesson:ZH-C159-yangtai}



\begin{quote}
Before the new one: say the Mandarin for a living room, then the Mandarin for a study at home.
\end{quote}

\subsection*{You'll want to know: yángtái}
\emph{yángtái} — \textquotedblleft{}balcony\textquotedblright{}.

\begin{quote}
\emph{wǒmen jiā de yángtái hěn dà} — Our balcony is big.
\end{quote}

Its characters come later; for now it is a word for the ear.

\begin{cousinweb}
\emph{yángtái} is built from \emph{yáng} 'the sunny side' and a word for a platform, both inherited from Middle Chinese, \emph{yang} and \emph{doj} (Guangyun, in Baxter's transcription). The dictionary writes the second another way and glosses it as raised square ground: a platform in the sun.
\end{cousinweb}

\begin{grammarlens}[title={What you've built}]
One more word for everyday things.
\end{grammarlens}

\subsection*{Your turn}
\begin{itemize}
\raggedright
  \item \emph{Say it:} \emph{yángtái}
  \item \emph{Say it:} \emph{yángtái}, once more
  \item \emph{Say it:} \emph{wǒmen jiā de yángtái hěn dà} — Our balcony is big.
  \item \emph{Recall:} say the Mandarin for a living room, then the Mandarin for a study at home, then say \emph{yángtái} again
\end{itemize}

\begin{tcolorbox}[breakable,colback=teal!4,colframe=teal!35!black,title={Before you move on}]
What does yángtái mean? (\textquotedblleft{}Balcony\textquotedblright{}.) Say it once more.
\end{tcolorbox}
\section[chuānghu]{chuānghu — window}

\label{lesson:ZH-C159-chuanghu}



\begin{quote}
Before the new one: say the Mandarin for a study at home, then the Mandarin for a balcony.
\end{quote}

\subsection*{You'll want to know: chuānghu}
\emph{chuānghu} — \textquotedblleft{}window\textquotedblright{}.

\begin{quote}
\emph{wòshì de chuānghu hěn dà} — The bedroom window is big.
\end{quote}

Its characters come later; for now it is a word for the ear.

\begin{cousinweb}
\emph{chuānghu} is built from \emph{chuāng} 'window' and \emph{hu} 'door', both inherited from Middle Chinese, \emph{tsrhaewng} and \emph{huX} (Guangyun, in Baxter's transcription). The dictionary writes the first another way and glosses it 'an opening right through'; it calls the second half of a door. In the word the second loses its tone.
\end{cousinweb}

\begin{grammarlens}[title={What you've built}]
One more word for everyday things.
\end{grammarlens}

\subsection*{Your turn}
\begin{itemize}
\raggedright
  \item \emph{Say it:} \emph{chuānghu}
  \item \emph{Say it:} \emph{chuānghu}, once more
  \item \emph{Say it:} \emph{wòshì de chuānghu hěn dà} — The bedroom window is big.
  \item \emph{Recall:} say the Mandarin for a study at home, then the Mandarin for a balcony, then say \emph{chuānghu} again
\end{itemize}

\begin{tcolorbox}[breakable,colback=teal!4,colframe=teal!35!black,title={Before you move on}]
What does chuānghu mean? (\textquotedblleft{}Window\textquotedblright{}.) Say it once more.
\end{tcolorbox}
